\section{Related Work}
\label{sec:related}

\paragraph{Minimum-edge-cut redistricting.}
Algorithmic redistricting has converged on edge-cut minimisation as a
mathematically clean, politically neutral criterion
\citep{chikina2017,deford2019,duchin2020}.
METIS multilevel graph partitioning \citep{karypis1998} is the standard solver:
it constructs a hierarchy of progressively coarsened graphs,
partitions the coarsest, and un-coarsens the partition back to the original.
\citet{karypis1998} show that METIS produces near-optimal solutions for
sparse irregular graphs in near-linear time.

\citet{b7paper} establish the seed-sweep convergence property: for states where the
minimum-edge-cut is unique (or nearly so), different METIS seeds converge
to the same partition, and the MEC plan is a property of the graph topology
rather than the seed.
This convergence property is a prerequisite for GeoSection: if the MEC at
each ratio is seed-independent, then the natural ratio selected by
isoperimetric normalisation is also seed-independent.

\paragraph{Compactness metrics.}
Single-district compactness is measured by Polsby-Popper
\citep{polsby1991} ($4\pi A/P^2$, where $A$ is district area and $P$ perimeter),
Reock \citep{reock1961} (ratio of district area to area of smallest enclosing circle),
and convex-hull ratio.
These metrics audit individual district shapes after the bisection tree
is complete.
Edge-cut minimisation is related but distinct: it minimises total boundary
length across all $k$ districts simultaneously, which is a globally stronger
criterion than minimising individual district perimeters.

The isoperimetric inequality that motivates our normalisation is classical:
for a convex planar region of area $A$ with boundary $P$, the minimum
perimeter enclosing fraction $f$ of the region satisfies
$P_{\min}(f) \propto \sqrt{f(1-f) \cdot A}$ \citep{osserman1978}.
We apply this principle to the edge-cut problem: the minimum-cost boundary
to separate $i$ out of $k$ population units scales as $\sqrt{i(k-i)/k^2}$
for uniform density, motivating the $\sqrt{\min(i,k-i)}$ normalisation.

\paragraph{Bisection ratio selection.}
Standard recursive bisection always uses the near-equal ratio \citep{karypis1998b}.
\citet{hendrickson1995} note that unequal partitioning can improve load balance
in parallel scientific computing, but do not study the normalised minimum-cut
selection criterion.
To our knowledge, GeoSection is the first redistricting algorithm to
systematically scan all feasible ratios at each bisection level with
isoperimetric normalisation.

\paragraph{Urban concentration and partisan effects.}
The geographic sorting of American voters --- Democrats concentrated in
dense urban cores, Republicans distributed across suburban and rural areas
--- is documented extensively \citep{rodden2019,chen2013}.
\citet{chen2013} show that geographic sorting produces Republican-leaning maps
under most compactness criteria without any intentional gerrymandering.
\citet{stephanopoulos2015} propose the efficiency gap as a measure of this
systematic bias.
\citet{mcghee2014} confirms that geography alone explains much of the
proportionality deficit in single-member district systems.

The caterpillar pathology is the redistricting algorithm's manifestation of
geographic sorting: because Democrats are more spatially concentrated,
minimum-edge-cut always finds their urban cores cheaper to isolate.
GeoSection addresses this by requiring the isolation to be confirmed
as genuinely compact per the isoperimetric standard.

\paragraph{GerryChain and Markov chain methods.}
\citet{deford2019} introduce ReCombiination (ReCom), a Markov chain over
redistricting plans that samples from the space of valid plans.
GerryChain \citep{metric2018} implements ReCom for practice.
These methods characterise the \emph{ensemble} of plans rather than
selecting a single optimal plan.
GeoSection is complementary: it selects a single deterministic plan
and provides a legal rationale (isoperimetric naturalness) that ensemble
methods do not directly supply.

\paragraph{Partisan fairness and legal standards.}
\citet{rucho2019} (\emph{Rucho v.\ Common Cause}) held that partisan
gerrymandering claims present a political question beyond the reach of
federal courts.
State courts have filled this gap:
\citet{lwv2018pa} (\emph{League of Women Voters v.\ Pennsylvania}) struck down
a congressional map as a partisan gerrymander under the Pennsylvania Constitution;
\citet{harper2022nc} (\emph{Harper v.\ Hall}) issued a similar ruling for
North Carolina.
Algorithmic redistricting that can articulate a politically neutral geometric
rationale is increasingly relevant to state-court review.
GeoSection's isoperimetric normalisation provides such a rationale:
the natural ratio is determined by geographic compactness, not partisan
composition.
